% Evidence-grading tables for the NS10657 supplemental material.
% Generated from code/analysis/evidence_grading.json (producer: 11_evidence_grading.ipynb).
% Three floats, one per output; labels tab:S4a/b/c. Included by supplemental_material.tex.

\begin{table*}[htbp]
\centering
\caption{Evidence grading, Output~A: terms that determine the absolute refractive index $n_{1762}$. Each source is graded \emph{quantified}, \emph{bounded} or \emph{unresolved}, matching panel~(b) of the evidence-grading figure of the main text. ``Magnitude'' is the assumed or measured size of the term; ``propagated effect'' states how it enters the refractive-index determination. Terms absorbed by the fitted intercept do not enter the environmental coefficients.}
\label{tab:S4a}
\footnotesize
\setlength{\tabcolsep}{4pt}
\begin{tabular}{@{}p{0.15\textwidth}p{0.095\textwidth}p{0.185\textwidth}p{0.095\textwidth}p{0.14\textwidth}p{0.20\textwidth}@{}}
\hline
\textbf{Source} & \textbf{Grade} & \textbf{Magnitude} & \textbf{Temporal behaviour} & \textbf{Evidence} & \textbf{Propagated effect} \\
\hline
Anchoring / interferometric offset \emph{(hard boundary; not correctable with this dataset)} & unresolved & $1.166\times10^{-6}$ (= 17.46 rad at 4.2 m; 4.16 rad/m) & static over the campaign; no independent monitor & nb 02; SM S1a; hard-boundary note & limits the absolute scale and baseline and blocks any ``Mathar absolutely correct'' claim; absorbed by the fitted intercept, so it does not enter the coefficients \\
K convention ($f_{780}/f_{1762}$ vs documented 1762.17/780.24) \emph{(convention / rounding, not a measurement)} & bounded & $dK/K = -5.17$ ppm $\to$ offset $5.2\times10^{-6}$, 4.4$\times$ the current $1.166\times10^{-6}$; 0.01 nm rounding alone allows $\pm$6.4 ppm & static multiplicative & nb 10 sec.~D; $f_{\rm rb}/f_{1762} = 2.2584857037$ & intercept only: $dK/K$ from $10^{-9}$ to $10^{-7}$ leaves $d(\alpha_H)_{\rm read}$ unchanged to four digits \\
780 nm frequency reference \emph{(Rb D2 saturated-absorption lock)} & quantified & $dn = 1.0\times10^{-9}$ ($d\nu \sim 384$ kHz at 384.230 THz) & static / long-term, wavemeter-monitored & SM S1a & enters single-epoch $n$ only; negligible for the coefficients \\
1762 nm quantum frequency reference & quantified & $dn < 10^{-12}$ & static & SM S1a & negligible \\
Interferometer non-linearity \emph{(declared negligible)} & bounded & not quantified & static & SM S1a & negligible \\
Residual scatter (incl.\ spatial/temporal mismatch) \emph{(statistical; the decomposition is not source identification)} & quantified & $1.84\times10^{-7}$; of which measurement noise $1.78\times10^{-7}$ and model nonlinearity $4.4\times10^{-8}$ & short-term scatter plus a slow component & SM S1a; nb 03 & sets single-epoch precision; the spatial/temporal mismatch inside it remains unresolved \\
Sensor-to-path spatial mismatch \emph{(possible contribution)} & unresolved & not separable; folded into the $1.84\times10^{-7}$ residual & short-term and slow & SM S1a & may contribute to the coefficient difference; must not be quoted as resolved \\
Wavelength convention (Mathar 1.762 vs 1.762174 $\mu$m; anchor 780.24 vs 780.246 nm) \emph{(documentation item)} & quantified & $d(dH) = -0.21\times10^{-12}$ and $+0.57\times10^{-12}$; anchor $dn = -0.037$ ppb & static & nb 10 sec.~D & $\leq 10^{-12}$ on the coefficients: no recomputation needed, only explicit documentation \\
\hline
\end{tabular}
\end{table*}

\begin{table*}[htbp]
\centering
\caption{Evidence grading, Output~B: terms that move the humidity coefficient difference $\Delta\alpha_H$, matching panel~(a) of the evidence-grading figure of the main text. ``Magnitude'' is the assumed or measured size of the term; ``propagated effect'' states how it enters the coefficient difference.}
\label{tab:S4b}
\footnotesize
\setlength{\tabcolsep}{4pt}
\begin{tabular}{@{}p{0.15\textwidth}p{0.095\textwidth}p{0.185\textwidth}p{0.095\textwidth}p{0.14\textwidth}p{0.20\textwidth}@{}}
\hline
\textbf{Source} & \textbf{Grade} & \textbf{Magnitude} & \textbf{Temporal behaviour} & \textbf{Evidence} & \textbf{Propagated effect} \\
\hline
Statistical uncertainty of the humidity difference \emph{(conditional on block length; not robust across the resampling procedures examined)} & quantified & $u_{\rm stat} = 1.30\times10^{-9}$ (SM nominal); $1.13\times10^{-9}$ at 0.56 h blocks; $3.69\times10^{-9}$ at 7.17 h blocks & block-length dependent; the paired-bootstrap SE grows from $1.13\times10^{-9}$ (0.56 h) to $5.38\times10^{-9}$ (24 h) & SM S1b; nb 06/07; \texttt{paired\_diff\_block\_sensitivity.json} & zero is excluded only at 0.56 h blocks; 7.17 h and 24 h blocks include zero \\
Humidity sensor gain (H channel) \emph{(consistent three-branch propagation, joint $(g,b)$ envelope)} & bounded & SM single-branch endpoint $3.99\times10^{-9}$ (corrected coordinate) vs joint envelope max $1.024\times10^{-9}$ ($u_{\rm syst} = 5.91\times10^{-10}$) & assumed constant multiplicative; BME280 $\pm$3 \%RH at $25\,^\circ$C, $\sim$4 \%RH over the campaign plus drift & nb 09; \texttt{joint\_gain\_offset\_envelope}; BME280 specification & moves $d(\alpha_H)_{\rm read}$ by $\leq 1.02\times10^{-9} = 24\%$ of the difference; gain is no longer the limiting systematic (the SM quote overestimates by 3.9$\times$) \\
Temperature and pressure sensor gain \emph{(same joint envelope)} & bounded & T: $1.03\times10^{-7} \to 1.484\times10^{-9}$; P: $1.23\times10^{-8} \to 2.324\times10^{-10}$ & assumed constant & nb 09 & cross-channel coupling is real and the response is not additive: the joint maximum exceeds the triangle sum of the channel-only maxima (H 1.024 vs $0.952\times10^{-9}$; linearised bound $0.915\times10^{-9}$), so a joint envelope is required \\
Sensor offset $b$ (centred parameterisation) \emph{(feasible envelope only; the two tutorial scenarios leave it at the endpoint gain)} & bounded & $|b| \leq \delta$; envelope vertex $b = -0.645$ (H) / $-0.125$ (T) / $-0.033$ (P); uncentred $b = (g-1)X_0 = 9.09$ \%RH at $g = 1.303$, 2.27$\times$ the spec bound & static & nb 09; task-card sec.~1 & the uncentred and the centred-$b=0$ hypotheses are different assumptions and both exceed the specification at the endpoint gain; only envelope vertices are legal sensitivity slices \\
Slow $1/R(t)$ component of the fringe-ratio anchor \emph{(shape not identifiable; bounded magnitude)} & unresolved & observed slow component 0.147--0.411 ppm; leverage $2.6$--$32.5\times10^{-9}$ per ppm; injection $+4.37\times10^{-9}$ (data-driven), up to $13.4\times10^{-9}$ (shape bound) & slow (ramp / seasonal); corr$(T,t) = 0.743$; the 111-day gap separates the two blocks & nb 10; \texttt{one\_over\_R\_term.json} & can account for the entire $4.33\times10^{-9}$ difference; it cannot be corrected with this dataset \\
Mathar validity-domain coverage \emph{(model-surrogate limitation, independent of calibration quality)} & unresolved & 91.7\% of rows above $25\,^\circ$C; in-domain $N = 12{,}134$ ($\sim$1.8 day equivalent), January plus 25 February rows; 28 intra-run gaps $> 2$ h (max $\sim$145 h) & static coverage property & nb 05; \texttt{campaign\_time\_coverage.json}; SM S0 & the in-domain comparison has limited discriminating power; the full-campaign comparison is an extrapolation and must be labelled as such \\
Block length for the coefficient uncertainties \emph{(decision pending, Task 4)} & unresolved & nominal $L = 156$ samples; 200 samples proposed & sets the statistical uncertainty, not a physical source & nb 07; todo v4 Task 4 & determines whether the difference excludes zero; no single nominal SE should be quoted before the decision \\
K as a reconstruction rather than a calibration \emph{(consistency check, not an independent validation)} & quantified & relative spread $3.95\times10^{-13}$; ptp/mean $2.98\times10^{-12}$ & no structure at this spread & nb 08; \texttt{chain\_propagation.calibrate\_K} & confirms that the toolset reproduces the archived pipeline; it does not validate the convention \\
\hline
\end{tabular}
\end{table*}

\begin{table*}[htbp]
\centering
\caption{Evidence grading, Output~C: refractivity change and slow phase compensation. ``Magnitude'' is the assumed or measured size of the term; ``propagated effect'' states how it enters the applied result.}
\label{tab:S4c}
\footnotesize
\setlength{\tabcolsep}{4pt}
\begin{tabular}{@{}p{0.15\textwidth}p{0.095\textwidth}p{0.185\textwidth}p{0.095\textwidth}p{0.14\textwidth}p{0.20\textwidth}@{}}
\hline
\textbf{Source} & \textbf{Grade} & \textbf{Magnitude} & \textbf{Temporal behaviour} & \textbf{Evidence} & \textbf{Propagated effect} \\
\hline
Offset-fair predictive improvement \emph{(over the observed interval and bandwidth)} & quantified & $I_{\rm amplitude} = 4.86\%$ [3.92, 5.95]; $I_{\rm variance} = 9.47\%$ [7.68, 11.54]; $\sigma_{\rm emp} = 1.8367\times10^{-7}$ vs $\sigma_{\rm Mathar} = 1.9304\times10^{-7}$ & campaign-scale & nb 02 & the surviving applied result; it must stay conditional on the observation interval and bandwidth \\
Frequency reference (780 nm / 1762 nm / K convention) \emph{(negative result for the design question Q2)} & quantified & 780 nm $dn = 1.0\times10^{-9}$; 1762 nm QFR $dn < 10^{-12}$; the K convention is intercept-only & static & SM S1a; nb 08/10 & a better frequency reference would not materially improve coefficient discrimination or slow compensation \\
Slow $1/R(t)$ drift \emph{(as Output~B row 5)} & unresolved & see Output~B row 5 & slow & nb 10 & limits predictive performance over intervals longer than the decorrelation time; it needs independent drift monitoring \\
Environmental coverage and observation schedule \emph{(campaign property, not a sensor limit)} & bounded & duty cycle 15.3\%; 111-day gap; in-domain $\sim$1.8 day equivalent & campaign-scale & nb 05; \texttt{campaign\_time\_coverage.json} & performance statements must be conditional on the schedule and the environmental range \\
Sensor gain and offset \emph{(as Output~B rows 2--4)} & bounded & joint envelope max $1.024\times10^{-9}$ (H) & static & nb 09 & small relative to the drift term for compensation; not the binding constraint \\
\hline
\end{tabular}
\end{table*}
